%% B-diagnostics.tex — Annex B: Diagnostic codes (normative)
\chapter{Diagnostic codes}
\label{annex:B-diagnostics}
\index{diagnostic codes}

\pnum
This annex is normative.  It lists every diagnostic code defined by this
specification, the clause that requires the diagnostic, the category
(constraint violation, warning, or informational), and a representative
message text.  A conforming implementation shall issue a diagnostic for
every constraint violation in the table below.

\pnum
Diagnostic codes are divided into the following ranges:

\begin{longtable}{@{}ll@{}}
\toprule
\textbf{Range} & \textbf{Domain} \\
\midrule
\endfirsthead
\toprule
\textbf{Range} & \textbf{Domain} \\
\midrule
\endhead
\bottomrule
\endlastfoot
\diagcode{E001}–\diagcode{E010}, \diagcode{E016}--\diagcode{E021} & Ownership and lifetime \\
\diagcode{E011}–\diagcode{E015}, \diagcode{E060} & Purity, safety, and verification \\
\diagcode{E082} & Bounded while verification \\
\diagcode{E083}–\diagcode{E087} & Array/slice safety and integer boundaries \\
\diagcode{E084}, \diagcode{E101}--\diagcode{E138} & Modules, attributes, generics, packing, ADT access, effects, verification \\
\diagcode{E100}                 & Miscellaneous \\
\diagcode{VRA}                  & Value range analysis \\
\end{longtable}

% -------------------------------------------------------------------------
\section*{Ownership and lifetime diagnostics (\diagcode{E001}–\diagcode{E010}, \diagcode{E016}, \diagcode{E020}–\diagcode{E021})}
% -------------------------------------------------------------------------

\begin{longtable}{@{}p{1.5cm}p{3.5cm}p{7.5cm}@{}}
\toprule
\textbf{Code} & \textbf{Clause} & \textbf{Message and trigger} \\
\midrule
\endfirsthead
\toprule
\textbf{Code} & \textbf{Clause} & \textbf{Message and trigger} \\
\midrule
\endhead
\bottomrule
\endlastfoot

\diagcode{E001}
  & \S\,\ref{sec:own.diag}
  & \textit{use of moved value '\texttt{x}'} ---
    Issued when an expression reads or passes a variable whose linear
    state is \textsc{Consumed}.
    \\[4pt]

\diagcode{E002}
  & \S\,\ref{sec:own.diag}
  & \textit{value '\texttt{x}' consumed twice} ---
    Issued when a second \lain{mov} expression is applied to a variable
    already in state \textsc{Consumed}. This subsumes the retired
    \texttt{E006} (\textit{consumed inside a loop}): consuming a value
    declared outside a loop from inside its body \emph{is} consuming it twice,
    and \diagcode{E002} names the violation where \texttt{E006} named
    only the syntactic setting it occurred in.
    \\[4pt]

\diagcode{E003}
  & \S\,\ref{sec:own.diag}
  & \textit{linear value '\texttt{x}' not consumed before end of scope} ---
    Issued when a \lain{mov}-mode variable exits scope in state
    \textsc{Unconsumed}.
    \\[4pt]

\diagcode{E004}
  & \S\,\ref{sec:own.diag}
  & \textit{cannot borrow '\texttt{x}' as mutable; it is already borrowed}
    --- Issued when a mutable borrow conflicts with an existing borrow
    (reads or write).
    \\[4pt]

\diagcode{E005}
  & \S\,\ref{sec:own.diag}
  & \textit{use of uninitialized variable '\texttt{x}'} ---
    Issued when a variable, struct field, or array element is read before
    it has been initialized (definite-assignment analysis).
    \\[4pt]



\diagcode{E007}
  & \S\,\ref{sec:own.diag}
  & \textit{moving linear value '\texttt{x}' requires an explicit
    \lain{mov}} --- Issued when a linear value is passed to an owned
    (\lain{mov}) parameter (or otherwise transferred) without the explicit
    \lain{mov} keyword: an implicit move.
    \\[4pt]

\diagcode{E010}
  & \S\,\ref{sec:mem.storage}
  & \textit{this reference would outlive the value it borrows} ---
    Issued when a reference whose referent has automatic storage duration
    would outlive it: returned (directly, or held by a struct or a sum,
    through any binding, assignment or call), stored into a place the
    caller owns, or passed to a function that keeps it there.
    \\[4pt]

\diagcode{E016}
  & \S\,\ref{sec:own.diag}
  & \textit{linear variable '\texttt{x}' is used inconsistently in
    the branches of \emph{stmt}} --- Issued when a linear value is
    consumed in some control-flow paths of an \lain{if}/\lain{case}
    but not in others, leaving an ambiguous state at the join point.
    Field-level inconsistency on linear struct fields is reported with
    the same code.
    \\[4pt]

\diagcode{E020}
  & \S\,\ref{sec:own.refbind}
  & \textit{cannot move out of a borrow} ---
    Issued when a \lain{mov}, or a whole-value copy of a linear value, takes
    its operand from a place the function does not own: a \lain{var} or
    shared parameter, a reference binding, a returned borrow, or a field or
    element of one.  The owner still holds the resource; moving it out would
    leave two owners.
    \\[4pt]

\diagcode{E021}
  & \S\,\ref{sec:own.diag}
  & \textit{assignment overwrites a linear value that still holds a resource}
    --- Issued when a linear place is assigned while it may still hold a
    resource, which would be lost.  Consume it first.  Through a borrow the
    old value can never have been consumed (\diagcode{E020}), so any such
    assignment is refused.
    \\[4pt]

\end{longtable}

% -------------------------------------------------------------------------
\section*{Purity, safety, and verification (\diagcode{E011}–\diagcode{E015})}
% -------------------------------------------------------------------------

\begin{longtable}{@{}p{1.5cm}p{3.5cm}p{7.5cm}@{}}
\toprule
\textbf{Code} & \textbf{Clause} & \textbf{Message and trigger} \\
\midrule
\endfirsthead
\toprule
\textbf{Code} & \textbf{Clause} & \textbf{Message and trigger} \\
\midrule
\endhead
\bottomrule
\endlastfoot

\diagcode{E011}
  & \S\,\ref{sec:fn.purity}
  & \textit{unacknowledged effect in '\texttt{f}'} ---
    Issued when a function's body has an effect its row does not name, while its
    bound is $\emptyset$ --- that is, while no \lain{effects} clause was written,
    or an empty one was.  A bound of $\emptyset$ is a function's default
    \emph{guarantee}, so violating it is this code; understating a \emph{non-empty}
    row is \diagcode{E130}.  The effects are \lain{io}, \lain{diverge},
    \lain{raises} and \lain{alloc}.  Also issued for a \lain{while} loop whose
    measure is neither given nor inferable, and for a direct self-recursion for
    which no measure can be inferred --- both sources of \lain{diverge}.  (A direct
    recursion \emph{with} a well-founded measure is permitted; a
    present-but-failing measure is \diagcode{E082}/\diagcode{E091}.)
    \\[4pt]

\diagcode{E012}
  & \S\,\ref{sec:interop.pointers}
  & \textit{type or constraint violation} ---
    Issued when a value does not have the type its context requires (an
    argument, a constructor field, a return, an operand, a \lain{case}
    scrutinee or pattern, \S\,\ref{sec:match.patterns}), when a value would
    have an opaque type or \lain{void}, when a refinement
    at a call site cannot be proven, and when a pointer-to-pointer or
    integer-to-pointer cast occurs outside an \lain{unsafe} block, and when a
    cast, inside \lain{unsafe} too, would view the bytes of a struct, or of a
    sum's payload, whose storage order the compiler chooses
    (\S\,\ref{sec:types.struct}; mark it \tok{[ordered]}).  The
    message names which; a test that expects \diagcode{E012} should say which
    it means, since the code alone does not.
    \\[4pt]

\diagcode{E013}
  & \S\,\ref{sec:basics.scopes}
  & \textit{redeclaration or shadowing of '\texttt{x}' is forbidden} ---
    Issued for a duplicate parameter name, a local that repeats or shadows
    an outer binding, a module-scope name defined twice (two \lain{extern}
    declarations of one function must agree), or a binding
    (variable, parameter, or loop variable) whose name is a builtin type
    (\lain{i32}, \lain{u8}, \lain{bool}, \ldots): a type name denotes a
    first-class type value and cannot also name a binding.
    \\[4pt]

\diagcode{E014}
  & \S\,\ref{sec:match.exhaustiveness}
  & \textit{non-exhaustive \lain{case}} ---
    Issued when the arms of a \lain{case} do not cover all values of the
    scrutinee type.  (An out-of-bounds \emph{index} is a different code,
    \diagcode{E085}.)
    \\[4pt]

\diagcode{E015}
  & \S\,\ref{sec:expr.arith}
  & \textit{a value that must be non-zero is not provably so} ---
    The divisor of a division or remainder, or the argument of \lain{@ctz} or \lain{@clz}
    (\S\,\ref{sec:types.vector}), when neither its type, its range nor a guard shows it is
    not zero.
    \\[4pt]

\end{longtable}

% -------------------------------------------------------------------------
\section*{Bounded while verification (\diagcode{E082})}
% -------------------------------------------------------------------------

\begin{longtable}{@{}p{1.5cm}p{3.5cm}p{7.5cm}@{}}
\toprule
\textbf{Code} & \textbf{Clause} & \textbf{Message and trigger} \\
\midrule
\endfirsthead
\toprule
\textbf{Code} & \textbf{Clause} & \textbf{Message and trigger} \\
\midrule
\endhead
\bottomrule
\endlastfoot

\diagcode{E082}
  & \S\,\ref{sec:stmt.while.bounded}, \S\,\ref{sec:fn.termination}
  & \textit{termination measure not proven well-founded} ---
    Issued when the implementation cannot verify that every looping-back path
    through the loop body strictly decreases the measure by at least 1
    (paths that \lain{break}/\lain{return} out are exempt), or that every
    self-call of a directly-recursive \lain{func} makes its measure strictly
    smaller, or that the measure is bounded below.  A measure the
    implementation cannot analyse at all is this code as well.

    \begin{note}
      \texttt{E080} (\textit{measure not proven non-negative}) and
      \texttt{E081} (\textit{cannot extract variables from measure}) were
      separate codes for two sub-cases of this one, issued by a second
      termination checker in the front end.  That checker is removed and the
      obligation has a single source, so the two codes are withdrawn: an
      absent or uninferable measure is \diagcode{E011}, a present one that
      cannot be shown well-founded is \diagcode{E082}, and nothing else is
      reported.  A conforming implementation shall not issue
      \texttt{E080} or \texttt{E081}.
    \end{note}
    \\[4pt]

\diagcode{E091}
  & \S\,\ref{sec:fn.termination}
  & \textit{malformed recursion measure} ---
    Issued when a directly-recursive \lain{func}'s \lain{decreasing} measure
    is not a parameter or a difference of two parameters.
    \\[4pt]

\end{longtable}

% -------------------------------------------------------------------------
\section*{Miscellaneous (\diagcode{E100})}
% -------------------------------------------------------------------------

\begin{longtable}{@{}p{1.5cm}p{3.5cm}p{7.5cm}@{}}
\toprule
\textbf{Code} & \textbf{Clause} & \textbf{Message and trigger} \\
\midrule
\endfirsthead
\toprule
\textbf{Code} & \textbf{Clause} & \textbf{Message and trigger} \\
\midrule
\endhead
\bottomrule
\endlastfoot

\diagcode{E100}
  & \S\,\ref{sec:05-lexical}, \S\,\ref{sec:07-types}
  & \textit{ill-formed token / semantic constraint violation} ---
    A catch-all code used for lexical and syntactic errors without a
    dedicated code, including: an unexpected token, using \lain{==} on enum
    types (\S\,\ref{sec:types.enum}), an empty or multi-character character
    literal, a keyword used as an identifier, and a zero-field array size.
    Also a construct this implementation does not model, refused where it
    appears rather than translated without it.
    \\[4pt]

\end{longtable}

% -------------------------------------------------------------------------
\section*{Array/slice safety and integer boundaries (\diagcode{E083}–\diagcode{E087})}
% -------------------------------------------------------------------------

\begin{longtable}{@{}p{1.5cm}p{3.5cm}p{7.5cm}@{}}
\toprule
\textbf{Code} & \textbf{Clause} & \textbf{Message and trigger} \\
\midrule
\endfirsthead
\toprule
\textbf{Code} & \textbf{Clause} & \textbf{Message and trigger} \\
\midrule
\endhead
\bottomrule
\endlastfoot

\diagcode{E083}
  & \S\,\ref{sec:types.struct}
  & \textit{linear field '\texttt{f}' missing \texttt{mov} annotation} ---
    Issued when a struct or enum variant field has a linear type but lacks
    the \lain{mov} keyword.  (Q-008 transitivity rule.)
    \\[4pt]

\diagcode{E085}
  & \S\,\ref{sec:expr.index}, \S\,\ref{sec:expr.subslice}
  & \textit{index or slice endpoint not proven in bounds} ---
    Issued when VRA cannot prove a non-negativity or upper-bound requirement
    for an array/slice index or sub-slice endpoint.
    \\[4pt]

\diagcode{E086}
  & \S\,\ref{sec:vra.boundary}
  & \textit{integer overflow at boundary: range exceeds target type} ---
    Issued when the VRA range of a source value provably exceeds the
    natural range of the target type at a type boundary (declaration,
    assignment, call argument, return, struct constructor), when a left shift may lose bits
    (\S\,\ref{sec:expr.bitwise}; \lain{<<\%} discards them on purpose), or when a module-scope constant's
    value, or an element of one, is not representable in its declared type
    (\S\,\ref{sec:decl.module-const}).
    \\[4pt]

\diagcode{E087}
  & \S\,\ref{sec:expr.subslice}, \S\,\ref{sec:vra.sized-slice-callsite}
  & \textit{sized slice constraint violated} ---
    Issued (1) when both endpoints of a sub-slice \lain{arr[a..b]} are
    integer literals and $a > b$; (2) at a call site where VRA can prove
    that the argument does not satisfy the formal parameter's sized slice
    constraint (\lain{T[k]}, \lain{T[>= k]}, \lain{T[> k]}, or
    \lain{T[ref.len]}).
    \\[4pt]

\end{longtable}

% -------------------------------------------------------------------------
\section*{Value range analysis (\diagcode{VRA})}
% -------------------------------------------------------------------------

\begin{longtable}{@{}p{1.5cm}p{3.5cm}p{7.5cm}@{}}
\toprule
\textbf{Code} & \textbf{Clause} & \textbf{Message and trigger} \\
\midrule
\endfirsthead
\toprule
\textbf{Code} & \textbf{Clause} & \textbf{Message and trigger} \\
\midrule
\endhead
\bottomrule
\endlastfoot

\diagcode{VRA}
  & \S\,\ref{sec:13-vra}
  & \textit{value range proof required} ---
    Issued when VRA cannot prove that an operation satisfies its
    required constraint and more specific code is not available.
    Appears as a diagnostic prefix in messages such as:
    ``\texttt{[VRA] cannot prove index in bounds}'' and
    ``\texttt{[VRA] return value may not satisfy constraint}''.
    \\[4pt]

\end{longtable}

\section*{Additional diagnostics (\diagcode{E008}, \diagcode{E009}, \diagcode{E017}--\diagcode{E125})}

The following diagnostics are emitted by the reference implementation and are
normative.  Each was verified against the compiler's diagnostic strings during
the \S\,\ref{annex:B-diagnostics} reconciliation pass.

\begin{longtable}{@{}ll@{}}
\toprule
\textbf{Code} & \textbf{Message and trigger} \\
\midrule
\endhead
\bottomrule
\endlastfoot
\diagcode{E008} & cannot move a value that is currently borrowed \\[2pt]
\diagcode{E009} & cannot mutate a field through a raw pointer in safe code --- use a \lain{var} borrow; or a write through a read-only raw pointer \lain{*T}, inside \lain{unsafe} too (\S\,\ref{sec:unsafe.enables}); or a write that reaches an immutable binding's or a module constant's storage through a reference, by any path (\S\,\ref{sec:stmt.assign}) \\[2pt]
\diagcode{E017} & a mutable borrow is spelled \lain{var}: passing to a \lain{var} parameter requires it at the call site, and a function returning \lain{var T} must \lain{return var} a place (or a call returning one) \\[2pt]
\diagcode{E018} & a function can reach its end without returning a value \\[2pt]
\diagcode{E019} & use of a partially-initialized value \\[2pt]
\diagcode{E060} & pointer dereference outside an \lain{unsafe} block \\[2pt]
\diagcode{E063} & a possibly-marker union value used without testing it first \\[2pt]
\diagcode{E064} & a \lain{|}-union cannot be niche-optimized (no spare bit pattern) \\[2pt]
\diagcode{E084} & access to a \lain{[private]} declaration from another module \\[2pt]
\diagcode{E101} & calling a \lain{proc} from a \lain{comptime} context (only \lain{func} is allowed) \\[2pt]
\diagcode{E102} & expected an attribute name after the opening bracket \\[2pt]
\diagcode{E103} & unknown attribute, or an attribute on a declaration it does not apply to (\tok{[ordered]} off a struct or a payload-carrying sum, \tok{[allocator]} on a function with a body) \\[2pt]
\diagcode{E104} & a type contains itself by value (infinite size) \\[2pt]
\diagcode{E105} & identifier is defined in another module --- add an \lain{import} \\[2pt]
\diagcode{E106} & use of an undeclared identifier, or a \lain{case} pattern that names no variant of the scrutinee's enum (\S\,\ref{sec:match.patterns}) \\[2pt]
\diagcode{E121} & a struct field invariant --- \lain{pos usize in text} (a valid index), \lain{pos usize <= src.len} or \lain{len usize <= cap} (a relation between fields), or a slice field's length bound \lain{src u8[<= 4294967295]} --- is not proven where the value is built or a field is written, or a \lain{var} reference is taken to a field it relates (a write through it would escape the check); also a \tok{[packed]} struct whose fields are not all \lain{iN}/\lain{uN} or exceed 64 bits, and a \lain{var} reference to, or the address of, one of its fields (bits have no address) \\[2pt]
\diagcode{E122} & a function pointer must be initialized with a matching function \\[2pt]
\diagcode{E123} & wrong number of arguments in a call through a function pointer \\[2pt]
\diagcode{E124} & generic instantiation: non-generic type, missing or uninferable type argument, a value argument that is not a constant, wrong arity \\[2pt]
\diagcode{E125} & direct ADT field access outside \lain{unsafe} --- destructure with \lain{case} \\[2pt]
\diagcode{E065} & a \lain{try} may propagate a marker the enclosing return union does not list \\[2pt]
\diagcode{E066} & \lain{try}/\lain{else} requires a \lain{T | markers} union value \\[2pt]
\diagcode{E067} & a checked operator (\lain{+?}, \lain{-?}, \lain{*?}, \lain{as?}) must be handled by \lain{else} \\[2pt]
\diagcode{E088} & a function returns a fixed-size array by value --- return a slice or write through a \lain{var} parameter \\[2pt]
\diagcode{E089} & a non-terminated value used where a sentinel-terminated type is required \\[2pt]
\diagcode{E090} & assignment between fixed-length strings of different lengths \\[2pt]
\diagcode{E126} & a mutable borrow (\lain{var}) stored in a struct field --- a field would outlive any loan the borrow checker can bound (a \emph{local} binding, \lain{var r = var p.x}, is admitted and held until its last use) \\[2pt]
\diagcode{E127} & call to an undeclared function \\[2pt]
\diagcode{E128} & the receiver has no such member, and no function of that name is in scope for UFCS \\[2pt]
\diagcode{E129} & \lain{assume} outside an \lain{unsafe} block --- use \lain{assert} to have the fact proven \\[2pt]
\diagcode{E130} & a declared, non-empty \lain{effects} row does not cover what the body does (an empty or absent row is \diagcode{E011}) \\[2pt]
\diagcode{E131} & an \lain{extern} declares a slice parameter --- a slice has no C representation; pass a pointer and a length \\[2pt]
\diagcode{E132} & a struct field's refinement is not one that is enforced --- a relation other than \lain{f CMP g} (g an integer field) or \lain{f CMP g.len} (g a slice or array field), a field's length bounded or sized by other than a constant, or a refinement after a slice type; any of them would be an unpaid assume \\[2pt]
\diagcode{E133} & a module-scope \lain{assert} whose operand is not a \lain{bool} constant expression (\S\,\ref{sec:decl.static-assert}) --- a run-time fact is asserted inside a function \\[2pt]
\diagcode{E134} & a module-scope assertion is false --- reported by the implementation, or, for one that measures a type with \lain{@sizeof}/\lain{@alignof}, by the C compiler through a \texttt{\_Static\_assert} carrying this code \\[2pt]
\diagcode{E135} & a \tok{[noreturn]} function can return: a \lain{return}, or the end of its body, is reachable (\S\,\ref{sec:decl.attributes}) \\[2pt]
\diagcode{E136} & a module constant computed at compile time cannot be: its initialiser reaches a function whose effect row is not empty, a function that may panic, or a call through a function pointer; its computation reads the constant itself, exceeds the step limit, or fails a check (\S\,\ref{sec:decl.module-const}) \\[2pt]
\diagcode{E137} & a compile-time constant is required and the expression is not one: a parameter refinement or sized-slice length the call site cannot evaluate (a call, a cast; \S\,\ref{sec:vra.sized-slice-callsite}), the length of a local array declared with an initializer, an array comprehension's bounds, a \lain{comptime if} condition \\[2pt]
\diagcode{E138} & control would leave a deferred statement: a \lain{return}, a \lain{break} or \lain{continue} aimed at a loop outside it, or a \lain{try} or \lain{else return} that leaves the function (\S\,\ref{sec:stmt.defer}) \\[2pt]
\end{longtable}
